\begin{pastexam}[label={px:ntu-108-dsaa-02}]{臺大資工所 108 學年度 第 2 題}
（問答）

(5\%) Convert the expression ``(a/(c*(b+d)))/(e-a)*c'' into a postfix form.
\end{pastexam}

\begin{pastsoln}
答案：$a\;c\;b\;d\;+\;*\;/\;e\;a\;-\;/\;c\;*$。

照資結書第~6.3.3~節的規則（虛擬碼~6.8）：運算元直接輸出；左括號推進堆疊；右括號把堆疊
彈到左括號為止；運算子先把堆疊頂端優先序\textbf{大於或等於}它的運算子彈出（同級左結合，
先出現的先算），再推進去；讀完後把堆疊全部彈出。 $*$、$/$ 同級，比 $+$、$-$ 高。 下表堆疊
由底寫到頂：
{\small
\begin{longtable}{clll}
\toprule
記號 & 動作 & 堆疊 & 輸出 \\
\midrule
\endfirsthead
\toprule
記號 & 動作 & 堆疊 & 輸出 \\
\midrule
\endhead
( & push & ( & \\
a & 輸出 & ( & a \\
/ & 頂端是 (，push & ( / & a \\
( & push & ( / ( & a \\
c & 輸出 & ( / ( & a c \\
$*$ & 頂端是 (，push & ( / ( $*$ & a c \\
( & push & ( / ( $*$ ( & a c \\
b & 輸出 & ( / ( $*$ ( & a c b \\
$+$ & 頂端是 (，push & ( / ( $*$ ( $+$ & a c b \\
d & 輸出 & ( / ( $*$ ( $+$ & a c b d \\
) & 彈出 $+$，丟掉 ( & ( / ( $*$ & a c b d $+$ \\
) & 彈出 $*$，丟掉 ( & ( / & a c b d $+$ $*$ \\
) & 彈出 /，丟掉 ( & （空） & a c b d $+$ $*$ / \\
/ & 堆疊空，push & / & a c b d $+$ $*$ / \\
( & push & / ( & a c b d $+$ $*$ / \\
e & 輸出 & / ( & a c b d $+$ $*$ / e \\
$-$ & 頂端是 (，push & / ( $-$ & a c b d $+$ $*$ / e \\
a & 輸出 & / ( $-$ & a c b d $+$ $*$ / e a \\
) & 彈出 $-$，丟掉 ( & / & a c b d $+$ $*$ / e a $-$ \\
$*$ & 頂端 / 同級，先彈出 /，再 push & $*$ & a c b d $+$ $*$ / e a $-$ / \\
c & 輸出 & $*$ & a c b d $+$ $*$ / e a $-$ / c \\
（結束） & 彈出全部 & （空） & a c b d $+$ $*$ / e a $-$ / c $*$ \\
\bottomrule
\end{longtable}
}

容易錯的是倒數第三步：$/$ 和 $*$ 同級，原式 $X / (e - a) * c$ 要讀成 $\bigl(X / (e - a)\bigr) * c$，
所以 $/$ 要先輸出，不能等到最後才和 $*$ 一起倒出來（那樣會得到 $\ldots e\,a\,-\,c\,*\,/$，
意思變成 $X / \bigl((e - a) * c\bigr)$）。

檢查：把結果用堆疊求值（資結書第~6.3.2~節）。 $a\,c\,b\,d\,+$ 得 $(b + d)$，接著 $*$ 得
$c * (b + d)$，$/$ 得 $a / (c * (b + d))$；$e\,a\,-$ 得 $(e - a)$，$/$ 得
$\bigl(a / (c * (b + d))\bigr) / (e - a)$；最後 $c\,*$，和原式一樣。
\end{pastsoln}
